\chapter{Poda con \emph{k-hop reachability}}
\label{cap:prune-khop}

\begin{intuicion}
NLF es un prune \emph{local}: chequea sólo el primer paso. Para
hipótesis multi-hop, querríamos descartar starts desde los que
\emph{no se puede llegar} al tipo destino final, sin importar qué
tan lejos. \emph{k-hop reachability} es exactamente eso: para cada
vértice \(v\) y cada tipo \(L\), un bit que dice ``existe un camino
de longitud \(\leq \texttt{MAX\_DEPTH}\) desde \(v\) a algún
vértice de tipo \(L\)''. Bit cero $\Rightarrow$ \(v\) no genera
match cuya hipótesis exija ese tipo. La extensión temporal
(\(\mathrm{min\_first}\), \(\mathrm{max\_first}\)) agrega el
chequeo de ventana: si la primera arista del camino más temprano
desde \(v\) cae fuera de la ventana del hunt, también se descarta.
Este capítulo formaliza el índice y la poda.
\end{intuicion}

\section{Definición}

\begin{definition}[k-hop reachability]
Sea \(\texttt{MAX\_DEPTH}\) una constante (en GraphHunter, \(4\)) y
\(L \in \Lc_V\) un tipo destino. Para \(v \in \V\):
\[
\bits[v][L] = 1 \iff \exists\,\text{walk}\;v = v_0, v_1, \dots, v_k\;
\text{con}\;k \leq \texttt{MAX\_DEPTH},\;\mathrm{type}(v_k) = L.
\]
\end{definition}

\noindent
La estructura concreta: \(\bits[v]\) es un \texttt{u16}, un bit por
\texttt{EntityType} (tag 0 a 8 ocupan 9 bits; los 7 restantes
quedan libres para extensiones futuras).

\begin{nota}
\(\bits\) es \emph{independiente del tipo de arista y del
timestamp}. Es una sobre-aproximación: ``\emph{podría} haber un
camino que cumpla''. Si bit es 1, el matcher no descarta; si es 0,
descarta seguro (porque ningún camino existe).
\end{nota}

\section{El prune}

\begin{algorithm}[H]
\caption{Prune k-hop al inicio del DFS}
\KwIn{starts candidatos \(\V_{L_0^V}\), hipótesis \(H\) con tipo destino final \(L_k^V\)}
\KwOut{starts filtrados}
\(L \leftarrow L_k^V.\mathrm{tag}\)\;
filtered \(\leftarrow \emptyset\)\;
\For{\(s \in \V_{L_0^V}\)}{
    \If{\(\bits[s] \,\&\, (1 \ll L)\)}{
        filtered.push(s)\;
    }
}
\Return filtered\;
\end{algorithm}

\noindent
Operación \emph{bitwise}: una sola instrucción de CPU
(\texttt{and} + comparación con cero). Para \(10^6\) starts, son
\(10^6\) operaciones de palabra, que es $\sim$1~ms.

\section{Soundness}

\begin{theorem}[k-hop prune sound]
Si la hipótesis \(H\) tiene longitud \(k \leq \texttt{MAX\_DEPTH}\)
y exige tipo destino \(L_k^V\), entonces:
\[
\bits[s][L_k^V.\mathrm{tag}] = 0 \implies s\;\text{no es origen de ningún match}.
\]
\end{theorem}

\begin{proof}
Por contradicción. Si existe match \(\sigma\) con \(\sigma(p_0) =
s\), entonces existe un walk \(s = v_0, v_1, \dots, v_k\) con
\(\mathrm{type}(v_k) = L_k^V\) y \(k \leq \texttt{MAX\_DEPTH}\)
(por hipótesis). Por definición de \(\bits\), eso implica
\(\bits[s][L_k^V] = 1\), contradicción. \(\square\)
\end{proof}

\subsection{Cuando \(k > \texttt{MAX\_DEPTH}\)}

Si la hipótesis es más larga que \(\texttt{MAX\_DEPTH}\), el prune
\emph{NO se aplica}: no podemos garantizar que un walk más largo no
exista aunque \(\bits = 0\). El motor en ese caso saltea el prune
y va al DFS directo.

En la práctica, \(\texttt{MAX\_DEPTH} = 4\) cubre el 95\% de las
hipótesis de threat hunting reales. Las hipótesis de longitud
\(\geq 5\) son raras (los analistas las escriben pero el matcher
les advierte que pueden ser lentas).

\section{Extensión temporal}

NLF y k-hop básico ignoran el tiempo. Para hipótesis con
\texttt{WITHIN}, querríamos descartar starts cuya \emph{primera
arista útil} caiga fuera de la ventana. La extensión D5/M2 agrega
dos matrices más:

\begin{definition}[\(\mathrm{min\_first}\), \(\mathrm{max\_first}\)]
\[
\mathrm{min\_first}[v][L] = \min \{\tau_0 : v \rightsquigarrow_{\leq \texttt{MAX\_DEPTH}} L\;\text{con primer edge}\;(\cdot,\cdot,\tau_0)\}
\]
\[
\mathrm{max\_first}[v][L] = \max \{\tau_0 : \dots\}
\]
\end{definition}

\noindent
Es decir: \(\mathrm{min\_first}[v][L]\) es el menor \emph{first edge
timestamp} entre todos los caminos de \(v\) a algún vértice de tipo
\(L\); \(\mathrm{max\_first}\) es el mayor.

\subsection{El prune temporal}

\begin{theorem}[Prune temporal sound]
Sea la hipótesis con ventana
\([t_{\mathrm{lo}}, t_{\mathrm{hi}})\). Si:
\[
\mathrm{min\_first}[v][L] \geq t_{\mathrm{hi}} \quad \text{or} \quad
\mathrm{max\_first}[v][L] < t_{\mathrm{lo}},
\]
entonces ningún walk de \(v\) al tipo \(L\) tiene su primer
\emph{edge} en la ventana.
\end{theorem}

\begin{proof}
Si \(\mathrm{min\_first}[v][L] \geq t_{\mathrm{hi}}\): todo walk
desde \(v\) tiene \(\tau_0 \geq t_{\mathrm{hi}}\), fuera de la
ventana por arriba. Si \(\mathrm{max\_first}[v][L] < t_{\mathrm{lo}}\):
todo walk tiene \(\tau_0 < t_{\mathrm{lo}}\), fuera por abajo.
En ambos casos, ningún walk válido existe. \(\square\)
\end{proof}

\subsection{Cuándo aplica}

El prune temporal se activa cuando la hipótesis tiene
\texttt{WITHIN}:

\begin{itemize}
    \item Hipótesis sin \texttt{WITHIN}: la ventana es
    \([-\infty, +\infty)\); ningún start se descarta por tiempo.
    \item Hipótesis con \texttt{WITHIN 30m} y un \emph{anchor time}
    explícito: el matcher chequea \(\mathrm{min/max\_first}\)
    versus la ventana.
\end{itemize}

\section{Storage layout: dense + sparse}

El layout naive sería: \(|\V| \times |\Lc_V|\) tabla densa para
cada uno de \(\bits\), \(\mathrm{min\_first}\),
\(\mathrm{max\_first}\). Para \(|\V| = 10^8\) vértices y
\(|\Lc_V| = 9\):

\begin{itemize}
    \item \(\bits\): \(10^8 \times 2\)~B (\texttt{u16}) = \(200\)~MB.
    \item \(\mathrm{min\_first}\): \(10^8 \times 9 \times 8\)~B
    (\texttt{i64}) = \(7.2\)~GB.
    \item \(\mathrm{max\_first}\): otros \(7.2\)~GB.
\end{itemize}

\noindent
Total: \(\sim 14.6\)~GB. Inviable.

La observación clave: \emph{muchos vértices no tienen aristas
salientes}. Para esos, \(\mathrm{min/max\_first}\) no aporta info
(no hay caminos). El layout \emph{sparse} guarda esas filas sólo
para vértices con \texttt{outdeg} > 0:

\begin{itemize}
    \item \(\bits\): denso \(|\V| \times 2\)~B.
    \item \(\mathrm{rank}\): \(|\V| \times 4\)~B (\texttt{u32}).
    \texttt{rank[v] = NO\_RANK} si \(\outdeg(v) = 0\); si no, indexa
    una fila en las matrices sparse.
    \item \(\mathrm{min\_first}\): sparse \(|\text{ranked}| \times 9
    \times 8\)~B.
    \item \(\mathrm{max\_first}\): igual.
\end{itemize}

\noindent
Para grafos típicos donde \(\sim 30\%\) de los vértices tienen
\texttt{outdeg} > 0 (\emph{ranked}):

\begin{itemize}
    \item \(\bits\): \(200\)~MB.
    \item \(\mathrm{rank}\): \(400\)~MB.
    \item \(\mathrm{min/max\_first}\) sparse: \(2 \times 30 \cdot 10^6
    \times 72\)~B = \(\sim 4.3\)~GB.
\end{itemize}

\noindent
Total: \(\sim 4.9\)~GB. Tres veces menos que el dense.

\begin{figure}[H]
\centering
\begin{tikzpicture}[
    node distance=0.4cm,
    every node/.style={font=\footnotesize, align=left},
    box/.style={draw=accent, thick, rounded corners=1mm,
                minimum height=1.4cm, minimum width=4.2cm,
                fill=accentlight}
]
\node[box] (bits) {\texttt{bits: Vec<u16>}\\\scriptsize dense, 2\,B/vertex\\\scriptsize 200\,MB @ 100M vértices};
\node[box, right=of bits] (rank) {\texttt{rank: Vec<u32>}\\\scriptsize dense, 4\,B/vertex\\\scriptsize 400\,MB @ 100M};
\node[box, below=of bits] (mn) {\texttt{min\_first: IndexBacking<[i64;9]>}\\\scriptsize sparse, 72\,B/ranked\\\scriptsize 2.16\,GB @ 30M ranked};
\node[box, right=of mn] (mx) {\texttt{max\_first: IndexBacking<[i64;9]>}\\\scriptsize sparse, 72\,B/ranked\\\scriptsize 2.16\,GB @ 30M ranked};
\end{tikzpicture}
\caption{\texttt{KHopReach} layout. Dense ($|\V|$) para \texttt{bits}
y \texttt{rank}; sparse (sólo vértices ranked) para
\texttt{min\_first} y \texttt{max\_first}. Total: $\sim$4.9~GB en
grafo de 100M vértices con 30\% \emph{outgoing-having}, vs $\sim$14.6~GB
del layout dense.}
\label{fig:khop-layout}
\end{figure}

\section{Construcción}

El \(\bits\) se construye con BFS multi-pasada (BFS reverso desde
todos los \emph{leaves} de cada tipo). El \(\mathrm{min/max\_first}\)
se calcula durante la misma pasada.

\begin{algorithm}[H]
\caption{KHopReach build (BFS multi-source)}
\KwIn{grafo \(G\)}
\KwOut{\texttt{KHopReach\{ bits, rank, min\_first, max\_first \}}}
\(\bits[v] \leftarrow 0\) para todo \(v\)\;
\For{\(L \in \Lc_V\)}{
    inicializar BFS desde todos los \(v\) con \(\mathrm{type}(v) = L\)\;
    \For{distance \(d = 0, 1, \dots, \texttt{MAX\_DEPTH}\)}{
        \For{cada vértice \(w\) en frontera \(d\)}{
            \(\bits[w] \mathrel{|}\!= (1 \ll L)\)\;
            actualizar \(\mathrm{min/max\_first}[w][L]\) con timestamps de las aristas reverse\;
        }
        avanzar frontera vía aristas reverse\;
    }
}
\Return\;
\end{algorithm}

\noindent
Costo: \(O(|\Lc_V| \cdot \texttt{MAX\_DEPTH} \cdot |\E|)\). Para
\(|\Lc_V| = 9\), \(\texttt{MAX\_DEPTH} = 4\), \(|\E| = 10^8\): \(\sim
3.6 \cdot 10^{10}\) operaciones, $\sim$1.5~min en hardware típico.
Lazy build (capítulo~\ref{cap:finalize}) lo difiere al primer
hunt.

\section{Resumen}

\begin{itemize}
    \item k-hop reachability: \(\bits[v][L] = 1\) sii existe walk
    de \(v\) al tipo \(L\) en \(\leq \texttt{MAX\_DEPTH}\) saltos.
    \item Prune sound: \(\bits[v][L] = 0 \Rightarrow v\) no genera
    match. No completo (sobre-aproxima).
    \item Extensión temporal: \(\mathrm{min/max\_first}[v][L]\)
    permite descartar starts cuya primera arista no cabe en la
    ventana.
    \item Layout dense + sparse para ahorrar memoria
    (\(\sim 3\times\) vs dense puro).
    \item Build \(O(|\Lc_V| \cdot \texttt{MAX\_DEPTH} \cdot |\E|)\),
    lazy en \texttt{OnceLock}.
\end{itemize}

\begin{ejercicios}
\begin{enumerate}
    \item \starred{} Para el grafo de la figura del capítulo
    \ref{cap:grafos-temporales}, calculá \(\bits[v]\) para cada
    vértice asumiendo \(\texttt{MAX\_DEPTH} = 2\).

    \item ¿Por qué la cota \(\texttt{MAX\_DEPTH} = 4\)? ¿Qué se
    rompe si la hipótesis es de longitud 5? ¿Y si fuera 6?

    \item \starred{} Para una hipótesis con WITHIN=30m y un
    anchor time específico, mostrá un ejemplo donde el prune
    temporal descarta un start aunque el prune k-hop básico no lo
    descarte.

    \item Calculá la memoria del layout sparse para un grafo de
    \(5 \cdot 10^7\) vértices con 50\% \emph{outgoing-having}.

    \item Diseñá un caso adversarial donde el prune k-hop tenga
    cero impacto (todos los starts pasan).

    \item ¿Qué cambia si en vez de un solo bit por tipo guardáramos
    \emph{cantidad} de caminos al tipo (no sólo si existe)?
    Discutí pros y contras.
\end{enumerate}
\end{ejercicios}
